\subsection{Sampling ordered orbit pairs}\label{p:sec:sampling}

We next approximate prescribed changes in finitely many coordinates by
positive orbit moves.  Both endpoint distributions must remain controlled:
this will allow the projection identity to pass from orbit pairs to the
couplings used in the product construction.  The scalar estimates of
Lemma~\ref{lem:scalar-sampling-estimates} apply coordinate by coordinate;
we prove here the growing-product estimate and the ordered weak limit.

The coordinate order is only Borel in the product-circle topology:
already in one circle, the ordered pairs $(1/2,1-1/n)$ converge to
the unordered pair $(1/2,0)$.  Thus the closed-relation hypothesis in
the classical coupling theorem \cite[Corollary~3.2]{Edwards1978} does
not apply directly.  The result below also requires approximating
measures supported on positive orbit pairs, with both marginal densities
controlled; those are additional conclusions proved by the sampling
construction.

For a finite set $I\subset\N$, define a finite Borel measure
$\lambda_I$ on $X\times X$ as follows.  In each coordinate $i\in I$
use the unnormalized triangular measure
\[
 \ind_{\{0\le x_i<y_i<1\}}\,dx_i\,dy_i,
\]
and in each other coordinate use the law of $(x_i,x_i)$ for uniform
$x_i\in[0,1)$.  Take the product over all coordinates.  Equivalently,
one can first normalize the finitely many triangular factors and
then multiply the resulting probability measure by $2^{-|I|}$.
In particular,
\[
 \lambda_I(X\times X)=2^{-|I|}.
\]
Its two marginal densities relative to $\mu$ are, respectively,
$\prod_{i\in I}(1-x_i)$ and $\prod_{i\in I}y_i$; they are bounded
by $1$.  The measure is carried by $x\le y$.

\begin{proposition}[Ordered orbit-pair sampling]\label{p:prop:sampling}
With probability one in the choice of $\theta$, the following holds
for every finite $I\subset\N$.  There are positive finite Borel
measures $\eta_{I,l}$ on $X\times X$ converging weakly to
$\lambda_I$, each carried by ordered pairs of the form
\[
 (x,y)=(x,\sigma_\theta^n x)
\]
with $n$ in a finite set of positive integers.  Both marginals of
$\eta_{I,l}$ have $L^2(\mu)$ densities, and their $L^2$ norms are
bounded uniformly in $l$.  The bound may depend on $I$ and on the
fixed realization of $\theta$.  Weak convergence here means
convergence against every continuous function on the compact space
$X\times X$.
\end{proposition}

\begin{proof}
We give the construction and its estimates before choosing the
probability-one event.  The notation $O(\cdot)$ in this proof has
absolute constants, uniform in the coordinate, the selected powers,
and the interval endpoints involved in a rectangle test.

We average many separated orbit powers so that both marginal densities
approach those of an ideal model with independent uniform phases.  The
coordinate weights impose triangular laws on $I$ and concentrate near
the diagonal elsewhere, while their marginal bounds remain uniform in
the growing product.  A final restriction to ordered tail coordinates
removes asymptotically negligible mass.

\paragraph{Separated powers and phase discrepancy}

Take $l$ sufficiently large that $m_l\ge200$, as required by the scalar
lemma. At stage $l$ write $m=m_l$, $d=d_l$, and set
\[
 Q=m^{(7/10)m},\qquad H=m^{300},\qquad R=N_lm^{1000}.
\]
For all sufficiently large $l$ we can choose, deterministically, a
set $\mathcal N_l$ of integers in $[Q/2,Q]$ such that
\begin{equation}\label{p:eq:sampling-separation}
 \begin{gathered}
 P_l:=|\mathcal N_l|\ge m^{0.69m},\\
 |vn+v'n'|>R\quad
 \left(\begin{array}{l}
 n,n'\in\mathcal N_l,\ n\ne n',\\
 v,v'\in\Z,\ |v|,|v'|\le H,\ (v,v')\ne(0,0)
 \end{array}\right).
 \end{gathered}
\end{equation}
Indeed, given a previously selected $n$, a fixed frequency pair with
$v'\ne0$ forbids an interval of possible $n'$ of length
$2R/|v'|\le2R$, hence at most $2R+2$ integers.  There are
$O(H^2)$ frequency pairs.  Since $N_l\le m^2$, each selection
therefore excludes at most $O(m^{1602})$ integers.  When $v'=0$
and $v\ne0$, the inequality is automatic because $n\ge Q/2>R$
eventually; likewise, a zero first frequency creates no obstruction.
The constraints are symmetric in the selected powers, since the
full set of frequency pairs is used.  Greedy selection from the
$\asymp Q$ available integers yields at least
$cQ/m^{1602}\ge m^{0.69m}$ choices for large $m$.

For a coordinate $i\in J_s$ with $s\le l$, the angle $\theta_i$ is
uniform on $[1/N_s,2/N_s]$.  Apply the separated-phase part of
Lemma~\ref{lem:scalar-sampling-estimates} with
\[
 N_0=N_s\le N_l=N\le m^2.
\]
The powers just selected satisfy its separation hypothesis.  Thus
$t_n=\{n\theta_i\}$ has a density $\varphi_{i,n}$ satisfying
\begin{equation}\label{p:eq:single-phase-density}
 \|\varphi_{i,n}-1\|_\infty\le N_s/n\le N_l/n.
\end{equation}
In particular, that density is at most two for large $l$.  For distinct
selected powers and functions $g,g'$ of bounded variation, the same
lemma gives
\begin{equation}\label{p:eq:bv-discrepancy}
 \begin{split}
 \left|\mathbb E_\theta g(t_n)g'(t_{n'})
       -\left(\int_0^1g\right)\left(\int_0^1g'\right)\right|
 &\le c m^{-80}
   (\|g\|_\infty+\operatorname{Var}g)
   (\|g'\|_\infty+\operatorname{Var}g').
 \end{split}
\end{equation}
No independence between $t_n$ and $t_{n'}$ is asserted.  The imported
estimate uses their Fourier coefficients and one-dimensional marginal
densities, and so applies even when their joint law is singular.
We shall also use the direct consequence of
\eqref{p:eq:single-phase-density}
\begin{equation}\label{p:eq:single-phase-test}
 \left|\mathbb E_\theta g(t_n)-\int_0^1g(t)\,dt\right|
 \le\frac{N_l}{n}\int_0^1|g(t)|\,dt.
\end{equation}

\paragraph{The weighted measures and both marginal orientations}

Fix a finite $I$ and restrict to sufficiently large $l$ that
$I\subset\{1,\ldots,d\}$.  Set
\[
 h=m^{-3/5},\qquad \Delta(s)=\min(s,1-s),\qquad \alpha=m^{-2}.
\]
For $0\le s,t<1$ define
\begin{equation}\label{p:eq:sampling-weights}
 w_i(s,t)=
 \begin{cases}
 \ind_{\{0<t<1-s\}},&i\in I,\\[1mm]
 \displaystyle\ind_{\{\Delta(s)\ge\alpha\}}
 \frac{\ind_{\{0<t<h\Delta(s)\}}}{h\Delta(s)},&i\notin I.
 \end{cases}
\end{equation}
The second line is defined to be zero if $\Delta(s)<\alpha$.
Let $\widetilde\eta_{I,l}$ be the average, over $n\in\mathcal N_l$,
of the push-forward under $x\mapsto(x,\sigma_\theta^n x)$ of
\begin{equation}\label{p:eq:untrimmed-weight}
 W_n(x)\,d\mu(x),\qquad
 W_n(x)=\prod_{i\le d}w_i(x_i,\{n\theta_i\}).
\end{equation}
Each weight is bounded for fixed $l$, so these are finite positive
measures.  They are carried by pairs satisfying $x_i<y_i$ for $i\le d$, since
$h\Delta(x_i)\le h(1-x_i)<1-x_i$ and $h<1$.

We will treat the marginals both without a test and after restriction
to a cylinder rectangle
\[
 \mathcal R=\{(x,y):x_i\in E_i^x,\ y_i\in E_i^y
                         \text{ for every }i\},
\]
where the $E_i^x,E_i^y$ are intervals with rational endpoints and
all but finitely many are the whole circle.  Take $l$ large enough
that all tested coordinates are at most $d$.  A missing test is
represented by using the whole circle for every interval.  For the
first and second marginal, respectively, define
\begin{equation}\label{p:eq:oriented-weights}
 \begin{split}
 u_i^{(1)}(s,t)
 &=w_i(s,t)\ind_{E_i^x}(s)\ind_{E_i^y}(\{s+t\}),\\
 u_i^{(2)}(s,t)
 &=w_i(\{s-t\},t)\ind_{E_i^x}(\{s-t\})\ind_{E_i^y}(s).
 \end{split}
\end{equation}
Changing variables $y=\sigma_\theta^n x$ preserves Haar measure.
It follows that the two tested marginal densities are exactly the
averages of the products of the corresponding $u_i^{(r)}$.
There is no factor from coordinates past $d$, even in the second
orientation, since their unweighted Haar law is translation invariant.

We check the extra factors introduced by triangular laws and interval
tests.  Write $u_i$ for either orientation and put
\[
 \bar u_i(s)=\int_0^1u_i(s,t)\,dt,\qquad
 A_i=\int_0^1\bar u_i(s)^2\,ds.
\]
Uniformly in $s$, the coordinate, and the test intervals, we have
\begin{equation}\label{p:eq:weight-variation}
 \|u_i(s,\cdot)\|_\infty+
       \operatorname{Var}_t u_i(s,\cdot)
 =O(m^{13/5})=O(m^{10}).
\end{equation}
For an untested triangular factor, the bound is $O(1)$: in the first
orientation its support in $t$ is $0<t<1-s$, and in the second it is
$0<t\le s$.  The possible value at $t=s$ affects no integral.
For the untested near-diagonal factors,
Lemma~\ref{lem:scalar-sampling-estimates} gives the bound
$O(m^{13/5})$.  To include the tests, in either orientation the
extra factor is a constant in $\{0,1\}$ times an interval indicator
translated or reflected modulo one.  Its supremum is at most one
and its variation is bounded by an absolute constant.  The inequality
\[
 \operatorname{Var}(ab)\le
 \|a\|_\infty\operatorname{Var}(b)
 +\|b\|_\infty\operatorname{Var}(a)
\]
therefore proves \eqref{p:eq:weight-variation}, including all jumps at
the test endpoints.

The moment bounds needed below are
\begin{equation}\label{p:eq:one-coordinate-moments}
 0\le\bar u_i(s)\le2,\qquad
 0\le A_i\le1+O(h^2),\qquad
 \int_0^1\int_0^1u_i(s,t)^2\,ds\,dt
       =O(h^{-1}\log m).
\end{equation}
For an untested triangular factor the two averages are exactly $1-s$
and $s$, respectively.  In either orientation its total integral and
squared integral are $1/2$, while $A_i=1/3$.  For an untested
near-diagonal factor, Lemma~\ref{lem:scalar-sampling-estimates} gives
\[
 \|\bar u_i\|_\infty\le2,\qquad
 A_i\le1+\frac{h^2}{(1-h)^2},\qquad
 \iint u_i(s,t)^2\,ds\,dt=\frac2h\log\frac1{2\alpha}.
\]
Since tests only decrease the nonnegative factors and their averages,
these statements imply \eqref{p:eq:one-coordinate-moments} for every
tested factor.  Their actual masses need not agree; in particular,
the near-diagonal integral identities are not being imposed on the
triangular factors or on the tested kernels.  The quadratic error
$O(h^2)$ is what keeps the product of the moment bounds controlled as
the number of coordinates grows.

\paragraph{The estimate in the growing product dimension}

For either fixed orientation and a fixed rectangle test or no test,
write
\[
 \begin{aligned}
 q_n(z)&=\prod_{i\le d}u_i(z_i,\{n\theta_i\}),&
 Z_l(z)&=\frac1{P_l}\sum_{n\in\mathcal N_l}q_n(z),\\
 Z_l^0(z)&=\prod_{i\le d}\bar u_i(z_i),&
 A&=\prod_{i\le d}A_i=\|Z_l^0\|_2^2.
 \end{aligned}
\]
We prove the estimate
\begin{equation}\label{p:eq:sampling-L2}
 \mathbb E_\theta\|Z_l-Z_l^0\|_{L^2(\mu)}^2
 \le O(dm^{-40})+\frac{(c h^{-1}\log m)^d}{P_l}.
\end{equation}
The independence used here is independence across coordinates
$\theta_i$.  In a single coordinate two phases may still be
dependent, as already accounted for in \eqref{p:eq:bv-discrepancy}.

For distinct $n,n'$, Tonelli's Theorem, product Haar measure, and
coordinate independence yield the exact factorization
\[
 \mathbb E_\theta\langle q_n,q_{n'}\rangle
 =\prod_{i\le d}\int_0^1
   \mathbb E_\theta\bigl[
    u_i(s,\{n\theta_i\})u_i(s,\{n'\theta_i\})\bigr]\,ds.
\]
By \eqref{p:eq:weight-variation} and \eqref{p:eq:bv-discrepancy}, its
$i$th factor is $A_i+e_i$ with $|e_i|\le c m^{-50}$.
Similarly,
\[
 \mathbb E_\theta\langle q_n,Z_l^0\rangle
 =\prod_{i\le d}\int_0^1
       \bar u_i(s)\mathbb E_\theta u_i(s,\{n\theta_i\})\,ds.
\]
Here \eqref{p:eq:single-phase-test} and $\bar u_i\le2$ bound the
error in each factor from $A_i$ by $4N_l/n=O(m^{-50})$.
Thus both sorts of factors have the same stated error bound.

The product errors do not require any lower bound on $A_i$.
Telescoping gives
\begin{equation}\label{p:eq:sampling-product-error}
 \left|\prod_{i\le d}(A_i+e_i)-\prod_{i\le d}A_i\right|
 \le cdm^{-50}(1+ch^2+cm^{-50})^{d-1}
 =O(dm^{-50}),
\end{equation}
because \eqref{p:eq:scale-consequences} gives $d\le1.01m$, and hence
$dh^2\le1.01m^{-1/5}$.  In particular,
$A\le(1+ch^2)^d=O(1)$.

For a self-pairing, use only the one-phase density bound at most $2$
and \eqref{p:eq:one-coordinate-moments}:
\[
 \mathbb E_\theta\|q_n\|_2^2
 =\prod_{i\le d}\int_0^1
            \mathbb E_\theta u_i(s,\{n\theta_i\})^2\,ds
 \le(c h^{-1}\log m)^d.
\]
To see the cancellation of all other main terms explicitly, expand
the squared norm.  The distinct-pair terms have total averaging
weight $1-1/P_l$ and main term $(1-1/P_l)A$.  The cross terms with
$Z_l^0$ have main term $-2A$, and the final squared norm of $Z_l^0$
is $A$.  Their sum is $-A/P_l\le0$.  By
\eqref{p:eq:sampling-product-error}, their remaining contribution
has absolute value $O(dm^{-50})$.  The self-pairings have total
weight $1/P_l$.  Dropping $-A/P_l$ proves the slightly weaker
bound \eqref{p:eq:sampling-L2}.

This estimate is summable in $l$.  Indeed,
\[
 \log_m (c h^{-1}\log m)^d
 \le1.01m\left(\frac35+\log_m(c\log m)\right)
 \le0.65m
\]
for large $m$, whereas $P_l\ge m^{0.69m}$.  Thus
\begin{equation}\label{p:eq:sampling-summable}
 \mathbb E_\theta\|Z_l-Z_l^0\|_2^2
 \le O(m^{-39})+m^{-0.04m},
\end{equation}
whose sum is finite by $m_l\ge2^l$.

\paragraph{One probability-one event and the weak limit}

For each fixed finite $I$, orientation, and rational cylinder
rectangle or absent test, Tonelli's Theorem applied to
\eqref{p:eq:sampling-summable} shows that
\[
 \sum_l\|Z_l-Z_l^0\|_2^2<\infty
 \quad\text{almost surely},
\]
where only the sufficiently large admissible stages are included.
In particular, $\|Z_l-Z_l^0\|_2\to0$.  Finite subsets of $\N$
form a countable family, as do finite cylinder rectangles with
rational endpoints, and there are only two orientations.  Intersect
these countably many probability-one events.  For the remainder
of the proof take any $\theta$ in this intersection.

Without a test, $Z_l$ is a marginal density of
$\widetilde\eta_{I,l}$.  Since $\|Z_l^0\|_2^2=A=O(1)$, the
$L^2$ convergence just proved bounds both marginal norms uniformly
in $l$.  We may omit finitely many initial stages; alternatively,
each of their densities is bounded for fixed $l$, so they can be
included by enlarging the bound.  The total masses of
$\widetilde\eta_{I,l}$ are consequently uniformly bounded.

For a fixed cylinder rectangle $\mathcal R$, using the first
orientation with its test shows that
\[
 \left|\widetilde\eta_{I,l}(\mathcal R)
       -\int_X Z_l^0\,d\mu\right|
 \le\|Z_l-Z_l^0\|_2\longrightarrow0.
\]
The ideal integral has a direct product interpretation on the first
$d$ coordinates.  For $i\in I$, integration of $w_i(x_i,t)$ in
$x_i,t$ gives exactly the triangular measure under $y_i=x_i+t$.
For $i\notin I$, first take uniform $x_i$, discard
$\Delta(x_i)<\alpha$, and, otherwise, take $t$ uniform in
$(0,h\Delta(x_i))$, setting $y_i=x_i+t$.  Before the discard this
is a probability law (the value at the two endpoint values of
$x_i$ may be defined arbitrarily), and
$|y_i-x_i|\le h/2$.  The total loss caused by all cutoffs is at most
\[
 2d\alpha=2dm^{-2}\longrightarrow0,
\]
by the union bound, with the finite triangular factors having mass
at most one.  Each fixed tested coordinate outside $I$ therefore
converges to its uniform diagonal law.  More explicitly, outside
an $h$-neighborhood of a tested interval's endpoints, membership
of $y_i$ in that interval agrees with membership of $x_i$; the
neighborhood has Lebesgue measure $O(h)$.  Since a fixed rectangle
tests only finitely many coordinates, all such errors vanish.
Untested coordinates contribute only their masses, and the total
cutoff estimate controls their product even as $d$ increases.
It follows that
\begin{equation}\label{p:eq:sampling-rectangle-limit}
 \widetilde\eta_{I,l}(\mathcal R)\longrightarrow
                  \lambda_I(\mathcal R).
\end{equation}

These limits imply weak convergence on $X\times X$.  Here are the
topological details.  The boundaries of the interval rectangles
have zero $\lambda_I$ measure, since every single-coordinate
marginal of $\lambda_I$ is absolutely continuous.  Every continuous
function on the compact product can be uniformly approximated by
a continuous function of finitely many coordinates: in a summable
product metric, replace the tail coordinates by fixed points and
use uniform continuity.  On the resulting finite torus, a sufficiently
fine rational grid gives a step function, a finite linear combination
of rational rectangle indicators, with arbitrarily small uniform
error.  Thus \eqref{p:eq:sampling-rectangle-limit} and the uniform
bound on total masses give convergence of integrals for every
continuous function.  In particular,
\[
 \widetilde\eta_{I,l}\ \Longrightarrow\ \lambda_I.
\]

\paragraph{Trimming the tail to impose order everywhere}

The untrimmed measures now have the required weak limit and both
marginal bounds.  It remains to impose order past the first $d$
coordinates.  Define
$\eta_{I,l}$ by restricting $\widetilde\eta_{I,l}$ to the Borel
set where $x_i\le y_i$ also for every $i>d$.  For a fixed selected
$n$, the weight $W_n$ depends only on the first $d$ coordinates,
so the tail $x$-coordinates remain independent uniform coordinates
under that summand, apart from multiplication by its total mass.
For a uniform tail coordinate, the probability that
$x_i>\{x_i+n\theta_i\}$ is exactly $\{n\theta_i\}\le n\theta_i$.
Consequently, uniformly over $n\in\mathcal N_l$, the relative mass
discarded from this summand is at most
\begin{equation}\label{p:eq:sampling-tail-trim}
 \begin{split}
 \sum_{i>d}n\theta_i
 &\le 2m^{(7/10)m}\sum_{s>l}\frac1{K_s}\\
 &\le\frac{4m^{(7/10)m}}{K_{l+1}}
 \le4m^{-m/20}\longrightarrow0.
 \end{split}
\end{equation}
The union bound is valid for the countable tail by increasing
finite unions.  More precisely, the discarded mass of a summand
equals its total mass times the Haar probability of its tail
failure event, because the weight and that event depend on disjoint
coordinates.  This statement also holds for a summand of mass zero.
Averaging over $n$ shows that the total discarded mass is bounded
by \eqref{p:eq:sampling-tail-trim} times the uniformly bounded
untrimmed mass.  Hence the difference of the two measures tends
to zero in total variation and
$\eta_{I,l}\Longrightarrow\lambda_I$.

Each trimmed marginal is a submeasure of the corresponding
untrimmed marginal.  Its Radon--Nikodym density is therefore
nonnegative and bounded almost everywhere by the untrimmed density,
so its $L^2$ norm cannot increase.  Trimming preserves the finite
set of positive orbit powers and now imposes order in every
coordinate.  All conclusions hold on the single probability-one
event already constructed, for every finite $I$, proving the
Proposition.
\end{proof}

We now fix one realization $\theta$ for which rational independence,
Proposition~\ref{p:prop:phases}, and Proposition~\ref{p:prop:sampling}
all hold, and henceforth write $\sigma=\sigma_\theta$.  The scales
and this rotation are fixed before any candidate invariant
projection is introduced.  For this fixed rotation, each sampling
sequence has actual weak convergence against every continuous
function; its validity is no longer restricted to the countable
tests used to obtain the probability-one event.  Thus later choices
of continuous tests or compact sets depending on a candidate
projection require no further intersection of events in the angle
probability space.
